\chapter*{Abstract}

The project entitled ``GG-Code: A Domain specific language for Generating and Editing G-code easier'' was developed by students from the Technical University of Moldova. This project report consists of the domain analysis, problem definition, proposed solution architecture, and system requirements, followed by conclusions and bibliography.

\textbf{Keywords: } CNC, G-code, Domain-Specific Language (DSL), Digital Manufacturing, Automation, Compiler Design.

This report presents the research, design, and initial implementation of a Domain-Specific Language (DSL) aimed at revolutionizing the workflow for Computer Numerical Control (CNC) and 3D printing systems. In the current digital manufacturing landscape, there exists a significant disconnect between high-level design intent and the low-level, verbose G-code instructions used to control machinery. Users are often forced to rely on opaque ``black box'' CAM software or resort to error-prone manual editing of cryptic machine codes.

The purpose of this project is to bridge this gap by introducing a human-readable, intermediate modelling language. This DSL allows manufacturing engineers, researchers, and hobbyists to define machine operations using programmable constructs such as variables, loops, and modular functions. The general objective is to develop a language specification and a compiler that translates these high-level descriptions into safe, optimized, and standard G-code.

The methodology employed involves a comprehensive domain analysis of current manufacturing toolchains, a Problem-Based Learning (PBL) assessment of user pain points, and a comparative analysis of existing solutions. The report details the formal grammar specification of the language and the implementation of a functional parser using ANTLR. This parser successfully processes complex constructs including motion control, geometric primitives, and algorithmic logic, validating the feasibility of the proposed solution.
